\begin{statementbox}
设双曲线
\[
H:\frac{x^2}{a^2}-\frac{y^2}{b^2}=1\quad(a,b>0),\qquad
c^2=a^2+b^2,\qquad e=\frac ca>1.
\]
两条渐近线是 $y=\pm\dfrac ba x$。
从\textbf{任意一个焦点} $F$ 出发，作直线 $\ell$ 平行于\textbf{其中一条}渐近线；
令 $\ell$ 与 $H$ 交于 $M$，与\textbf{另一条}渐近线交于 $N$。
则 $F,M,N$ 三点按 $F\!\to\!M\!\to\!N$ 的顺序排列。

\textbf{核心定比：}
\[
\boxed{\frac{FM}{MN}=\frac{b^2}{a^2}=e^2-1
      =\tan^2\theta},\qquad \tan\theta=\frac ba.
\]
其中 $\theta\in(0,\pi/2)$ 是渐近线与 $x$ 轴的锐角。

\begin{center}\begin{tikzpicture}[font=\small,>=stealth]
\draw[gray!45,rounded corners=3pt] (0,0) rectangle (14.2,10);
\node at (7.1,9.35) {先认三点：过焦点作平行线，沿线依次遇到 M、N};

\begin{scope}[shift={(4.4,4.8)},scale=.7]
\draw[gray,->] (-5.2,0)--(5.8,0) node[right] {$x$};
\draw[gray,->] (0,-4.6)--(0,4.6) node[above] {$y$};
\draw[violet,dashed,thick] (-3.2,-64/15)--(3.2,64/15);
\draw[violet,dash dot,thick] (-3.2,64/15)--(3.2,-64/15);

\draw[teal!75!black,thick,samples=90,domain=-4.2:4.2,variable=\v] plot ({3*sqrt(1+\v*\v/16)},{\v});

\draw[teal!75!black,thick,samples=90,domain=-4.2:4.2,variable=\v] plot ({-3*sqrt(1+\v*\v/16)},{\v});

\coordinate (F) at (5,0);\coordinate (M) at (3.4,-32/15);\coordinate (N) at (2.5,-10/3);
\draw[blue!75!black,very thick] (F)--(M);
\draw[orange!85!black,very thick] (M)--(N);
\foreach \q in {F,M,N}\fill (\q) circle(2.8pt);
\node[above right=4pt] at (F) {$F$};\node[right=5pt] at (M) {$M$};\node[below left=4pt] at (N) {$N$};

\node[below left] at (0,0) {$O$};
\node[violet,fill=white,inner sep=2pt] at (-2,3.6) {$y=-kx$};
\node[violet,fill=white,inner sep=2pt] at (-2,-3.6) {$y=kx$};
\draw[blue!75!black,->] (5.7,-1)--(4.35,-.87);
\node[blue!75!black,anchor=west] at (5.8,-1) {$FM$};
\draw[orange!85!black,->] (4.1,-3.6)--(2.95,-2.73);
\node[orange!85!black,anchor=west] at (4.2,-3.6) {$MN$};
\end{scope}
\node[anchor=west,align=left] at (9.5,6.9) {$k=b/a$\\[10pt]过 $F$ 的线 $\ell$\\[8pt]平行 $y=kx$};
\node[anchor=west,align=left] at (9.5,4.4) {$M$ 在双曲线上\\[10pt]$N$ 在另一条\\[6pt]渐近线上};
\node at (7.1,1.1) {本图 $a=3,b=4,c=5$：蓝段 $8/3$，橙段 $3/2$，比为 $16/9$。};
\node at (7.1,.5) {两段方向相同；$FN=FM+MN$。};

\end{tikzpicture}
\end{center}
\end{statementbox}
\clearpage
\begin{statementbox}
若题设为向量关系
\[
\overrightarrow{FM}=\lambda\overrightarrow{MN}\quad(\lambda>0),
\]
则由于两向量同向，$\lambda=FM/MN$，所以
\[
\boxed{e=\sqrt{1+\lambda}}.
\]


\textbf{强化结论：三段定长。}
同一构型中，还能求出三段的\textbf{绝对长度}：
\[
\boxed{FM=\frac{b^2}{2a}},\qquad
\boxed{MN=\frac a2},\qquad
\boxed{FN=\frac{c^2}{2a}}.
\]
特别地，$MN=a/2$ 与 $b$ 无关；但\textbf{整条图形的形状和位置}仍随 $b$ 改变。

\textbf{对称性：}取左焦点或右焦点，平行上升或下降的渐近线，
只会将图像反射；以上比例与长度均不变。

\textbf{需要保留的条件：}
\begin{itemize}
\item 双曲线实轴沿 $x$ 轴，且 $a,b>0$；
\item $\ell$ \textbf{过焦点且平行一条渐近线}；
\item $N$ 是 $\ell$ 与\textbf{另一条}渐近线的交点；
\item 向量等式中 $\lambda>0$ 对应 $F-M-N$，不能把 $FM/MN$ 误写成倒数。
\end{itemize}

\textbf{认知主线：}
\[
\boxed{\text{焦点 + 平行渐近线}}
\ \Longrightarrow\
\boxed{\text{一次交点}}
\ \Longrightarrow\
\boxed{\frac{FM}{MN}=\left(\frac ba\right)^2=e^2-1}.
\]
\end{statementbox}
